\chapter{Compact \(a_2/a_4\) spectral-moment gate}

Внешний red-team предложил использовать compact heat-kernel coefficients
для фиксации quadratic и quartic terms одним spectral action. Выполним
это вычисление в той же Laplace-конвенции, что использовалась для
\(\kappa\) и \(g\).

\section{Curved product Dirac convention}

На четырёхмерном base-\(K\) оператор имеет вид
\begin{equation}
 \mathcal D_A
 =
 i\gamma^\mu(\nabla_\mu+A_\mu)
 +\gamma^5\Phi.
\end{equation}
В Laplace-записи
\begin{equation}
 \mathcal D_A^2
 =-(\nabla^2+E)
\end{equation}
endomorphism содержит
\begin{equation}
 E
 =
 -\frac14R\mathbf1
 -\Phi^2
 -i\gamma^\mu\gamma^5D_\mu\Phi
 -\frac12\gamma^{\mu\nu}F_{\mu\nu}.
\end{equation}
Для closed manifold relevant coefficients равны
\begin{align}
 a_2
 &=
 \frac1{(4\pi)^2}
 \int_K
 \Tr\left(E+\frac16R\right),
 \\
 a_4
 &=
 \frac1{(4\pi)^2}
 \int_K
 \Tr\left(
 \frac12E^2
 +\frac16RE
 +\frac1{12}\Omega^2
 +\text{pure curvature}
 \right),
\end{align}
где total derivatives опущены.

\section{\(\Phi\)-часть коэффициента \(a_2\)}

Подстановка \(E\) даёт
\begin{equation}
 E+\frac16R
 =
 -\Phi^2-\frac1{12}R
 +\text{Clifford-odd terms}.
\end{equation}
После spin trace
\begin{equation}
 a_2\big|_{\Phi}
 =
 \frac1{(4\pi)^2}
 \int_K
 \left[
 -4\,\Tr_{\mathrm{orb}}\Phi^2
 \right].
\end{equation}
Коэффициент \(a_0\) содержит только
\(\Tr\mathbf1\) и не зависит от \(\Phi\).

\section{\(\Phi\)-часть коэффициента \(a_4\)}

Quartic term приходит из \(\frac12E^2\):
\begin{equation}
 a_4\big|_{\Phi^4}
 =
 \frac1{(4\pi)^2}
 \int_K
 2\,\Tr_{\mathrm{orb}}\Phi^4.
\end{equation}
Derivative term имеет тот же spin-traced coefficient:
\begin{equation}
 a_4\big|_{\mathrm{kin}}
 =
 \frac1{(4\pi)^2}
 \int_K
 2\,\Tr_{\mathrm{orb}}(D_\mu\Phi D^\mu\Phi).
\end{equation}

Curvature--scalar term получает два вклада. Из
\(\frac16RE\) приходит
\(-R\Phi^2/6\) до spin trace, а cross-term
\(\frac12E^2\) между \(-R/4\) и \(-\Phi^2\) даёт
\(+R\Phi^2/4\). Сумма равна
\begin{equation}
 -\frac16R\Phi^2+\frac14R\Phi^2
 =\frac1{12}R\Phi^2.
\end{equation}
После spin trace
\begin{equation}
 a_4\big|_{R\Phi^2}
 =
 \frac1{(4\pi)^2}
 \int_K
 \frac13R\,\Tr_{\mathrm{orb}}\Phi^2.
\end{equation}

\begin{theorem}[Compact scalar heat-kernel ledger]
В выбранной convention \(\Phi\)-dependent coefficients равны
\begin{align}
 a_2\big|_\Phi
 &\propto-4\Tr_{\mathrm{orb}}\Phi^2,
 \\
 a_4\big|_\Phi
 &\propto
 2\Tr_{\mathrm{orb}}(D\Phi)^2
 +2\Tr_{\mathrm{orb}}\Phi^4
 +\frac13R\Tr_{\mathrm{orb}}\Phi^2.
\end{align}
Тем самым \(a_2\) фиксирует quadratic bare term, а \(a_4\) одновременно
фиксирует kinetic, quartic и conformal curvature coupling.
\end{theorem}

\section{Нормированный bare potential}

Четырёхмерный spectral action содержит
\begin{equation}
 S_{\mathrm{spec}}
 \supset
 f_2\Lambda^2a_2+f_0a_4.
\end{equation}
После деления на общий scalar kinetic coefficient \(2f_0\) получаем
\begin{align}
 \mathcal L_{\Phi}
 &=
 \Tr_{\mathrm{orb}}(D_\mu\Phi D^\mu\Phi)
 +\Tr_{\mathrm{orb}}\Phi^4
 \\
 &\quad
 +\left(
 -2\frac{f_2}{f_0}\Lambda^2
 +\frac16R
 \right)
 \Tr_{\mathrm{orb}}\Phi^2.
\end{align}
Сопоставление с
\begin{equation}
 \Tr_{\mathrm{orb}}
 \left(\Phi^2-\chi^2\mathbf1\right)^2
\end{equation}
даёт
\begin{equation}
 \boxed{
 \chi^2
 =
 \frac{f_2}{f_0}\Lambda^2
 -\frac1{12}R.}
\end{equation}

Для круглого
\(K=\mathbb{RP}^3_R\times S^1_R\) с
\begin{equation}
 R_K=\frac6{R^2}
\end{equation}
получаем dimensionless relation
\begin{equation}
 \boxed{
 (\chi R)^2
 =
 \frac{f_2}{f_0}(\Lambda R)^2
 -\frac12.}
\end{equation}
Ненулевой flattening vacuum существует только если
\begin{equation}
 \frac{f_2}{f_0}(\Lambda R)^2>\frac12.
\end{equation}

\section{Что закрыто}

\begin{enumerate}
 \item bare quartic/kinetic ratio фиксирован и равен единице;
 \item curvature coupling фиксирован и не является новым coefficient;
 \item quadratic term определяется теми же spectral data
 \(f_2,\Lambda,f_0\);
 \item прежняя свободная \(\lambda_2\) на bare matching scale заменена
 одной spectral combination.
\end{enumerate}

\begin{proposition}[Tree-level spectral boundary pass]
Compact \(a_2/a_4\) calculation действительно задаёт единое tree-level
boundary condition для scalar potential. Независимое ручное назначение
\(\lambda_2\) и \(\lambda_4\) на matching scale больше не требуется.
\end{proposition}

\section{Что не закрыто}

Формула для \(\chi R\) содержит
\begin{equation}
 \zeta_{\mathrm{mom}}
 :=
 \frac{f_2}{f_0}(\Lambda R)^2.
\end{equation}
Ни topology, ни finite algebra, ни charge trace пока не фиксируют
\(\zeta_{\mathrm{mom}}\). Выбор spectral function определяет
\(f_2/f_0\), а \(\Lambda R\) задаёт cutoff относительно геометрического
радиуса.

Кроме того, one-loop renormalized quartic допускает finite matching shift,
если не постулировано, что bare spectral boundary condition является
полной renormalization condition при \(\mu=\Lambda\).

\begin{nogocriterion}[Spectral-moment closure rule]
Scale gate считается закрытым только если:
\begin{enumerate}
 \item spectral function или её moment ratio \(f_2/f_0\) выведены;
 \item relation \(\Lambda R\) фиксирована геометрически или динамически;
 \item finite quantum matching shift при \(\mu=\Lambda\) доказанно
 отсутствует либо вычислен.
\end{enumerate}
Иначе \(\zeta_{\mathrm{mom}}\) является одной непрерывной scale-setting
комбинацией.
\end{nogocriterion}

\begin{theorem}[Compact moment verdict]
Предложенный external red-team маршрут частично проходит: он закрывает
двухпараметрическую свободу bare \(\lambda_2,\lambda_4\) до одной
spectral-moment combination. Но parameter-free absolute scale не получен.
После одного train scale relation становится проверяемой boundary
condition для \(f_2/f_0\) и \(\Lambda R\).
\end{theorem}

\begin{protocol}[Следующий гейт]
Необходимо проверить минимальные заранее объявленные spectral functions
и вычислить их \(f_2/f_0\). Для каждой функции следует:
\begin{enumerate}
 \item зафиксировать normalization \(f_0\);
 \item вычислить \(f_2/f_0\);
 \item проверить vacuum inequality;
 \item определить, остаётся ли \(\Lambda R\) свободным.
\end{enumerate}
Сравнение с particle masses до этого запрещено.
\end{protocol}

Машинный аудит сохранён в
\path{s2t_v3_compact_a2_a4_moment_results.json}.